\section{Factorization over abelian base fields}

Let $\Phi_k$ be the $k$-th cyclotomic polynomial, let $\varphi$ be
Euler's totient function, and define
\[
 \Gamma_k(T)=(-1)^{\varphi(k)}T^{2\varphi(k)}\Phi_k(g(T)).
\]
For $k\ge2$, the identity \eqref{eq:draft-g-F} gives
\begin{equation}\label{eq:Gamma-fiber-product}
 \Gamma_k(T)=
 \prod_{\substack{\zeta^k=1\\ \zeta\ \mathrm{primitive}}}F_\zeta(T).
\end{equation}
By Theorem \ref{thm:abelian-stability}, for every nontrivial root of
unity $\zeta$, the polynomial $F_\zeta$ is irreducible over
$\mathbf Q(\zeta)^{\mathrm{ab}}$, and for every root $\alpha$ one has
\[
 \mathbf Q(\zeta,\alpha)\cap
 \mathbf Q(\zeta)^{\mathrm{ab}}=\mathbf Q(\zeta).
\]
Fix a primitive $k$-th root $\zeta_k$, put
$K_k=\mathbf Q(\zeta_k)$, and let $B/\mathbf Q$ be a finite abelian
extension.  Set
\[
 D=B\cap K_k.
\]
Because both $B/\mathbf Q$ and $K_k/\mathbf Q$ are Galois, restriction
gives
\[
 \operatorname{im}\!\left(\operatorname{Gal}(\overline{\mathbf Q}/B)
 \longrightarrow\operatorname{Gal}(K_k/\mathbf Q)\right)
 =\operatorname{Gal}(K_k/D).
\]
The orbits on the primitive $k$-th roots therefore have cardinality
$[K_k:D]$.  Thus
\[
 \Phi_k(X)=\Psi_1(X)\cdots\Psi_r(X)\quad\text{in }B[X],
 \qquad
 r=[D:\mathbf Q],\quad \deg\Psi_i=[K_k:D].
\]
Here the $\Psi_i$ are the distinct monic irreducible factors over $B$.
For each $i$, define
\[
 \Gamma_{\Psi_i}(T)
 =(-1)^{\deg\Psi_i}T^{2\deg\Psi_i}\Psi_i(g(T)).
\]
The fixed-degree identity $T^2(g(T)-\zeta)=-F_\zeta(T)$ gives
\[
 \Gamma_{\Psi_i}(T)=\prod_{\Psi_i(\zeta)=0}F_\zeta(T)
 \quad\text{and}\quad
 \Gamma_k=\prod_{i=1}^r\Gamma_{\Psi_i}.
\]

Fix a root $\zeta$ of $\Psi_i$ and a root $\alpha$ of $F_\zeta$.
Since $F_\zeta(0)=-16$, one has $\alpha\ne0$, so the Laurent expression
$g(\alpha)$ is defined.
The extension $BK_k/K_k$ is abelian, hence is contained in
$K_k^{\mathrm{ab}}$.  Therefore $F_\zeta$ remains irreducible over
$BK_k$, and
\[
 [BK_k(\alpha):BK_k]=6.
\]
On the other hand, the relation $\zeta=g(\alpha)$ shows that
$BK_k\subset B(\alpha)$, and in fact $B(\alpha)=BK_k(\alpha)$.  Hence
\[
 [B(\alpha):B]
 =6[BK_k:B]
 =6[K_k:D].
\]
This is the degree of $\Gamma_{\Psi_i}$.  Since
$\Gamma_{\Psi_i}\in B[T]$ is monic and vanishes at $\alpha$, it is the
minimal polynomial of $\alpha$ over $B$ and is irreducible.  These
irreducible polynomials are pairwise distinct: if an element $\alpha$
were a common root of $\Gamma_{\Psi_i}$ and $\Gamma_{\Psi_j}$, then
$g(\alpha)$ would be a common root of the coprime polynomials
$\Psi_i$ and $\Psi_j$.  This proves the following theorem.

\begin{theorem}\label{thm:abelian-factorization}
For $k\ge2$ and every finite abelian extension $B/\mathbf Q$, the
polynomial $\Gamma_k$ has exactly
\[
 [B\cap K_k:\mathbf Q]
\]
irreducible factors over $B$, each of degree
\[
 6[K_k:B\cap K_k].
\]
In particular,
\[
 \Gamma_k\text{ is irreducible over }B
 \quad\Longleftrightarrow\quad
 B\cap K_k=\mathbf Q.
\]
All splitting after an abelian base change is therefore caused by the
cyclotomic factor $\Phi_k$; no sextic fiber splits further.
\end{theorem}

\begin{corollary}[Factorization over the maximal abelian extension]
\label{cor:Qab-factorization}
For every $k\ge2$, the complete irreducible factorization of $\Gamma_k$
over $\mathbf Q^{\mathrm{ab}}$ is
\begin{equation}\label{eq:Qab-complete-factorization}
 \Gamma_k(T)=
 \prod_{\substack{\zeta^k=1\\ \zeta\ \mathrm{primitive}}}F_\zeta(T).
\end{equation}
It has $\varphi(k)$ factors, all of degree six.
\end{corollary}

\begin{proof}
Since $K_k\subset\mathbf Q^{\mathrm{ab}}$ and
$\mathbf Q^{\mathrm{ab}}/K_k$ is abelian, it is contained in
$K_k^{\mathrm{ab}}$.  Each $F_\zeta$ is consequently irreducible over
$\mathbf Q^{\mathrm{ab}}$.  Distinct primitive roots give distinct
polynomials, as their $T^2$ coefficients are $\zeta-5$.  Equation
\eqref{eq:Gamma-fiber-product} now proves the assertion.
\end{proof}

\section{Transfer to the period-4 delta factors}

We use the standard delta-factor notation
\cite{Morton1996,MortonVivaldi1995,KoizumiMurakamiSanoTakehira2025}.
Write $C=4c$ for the scaled quadratic-family parameter and put
\[
 \widetilde\Delta_{4k,4}(C)=\Delta_{4k,4}(C/4).
\]
For a polynomial $q$, let $Z(q)$ denote its set of roots in
$\overline{\mathbf Q}$, without multiplicity.  The period-$4$
normalization has parameter coordinate
\[
 h(T)=-T^2-3-\frac4T.
\]
The established normalization and separability results
\cite[Proposition 2.3, Equation (2.2), and Remark 2.10]
{KoizumiMurakamiSanoTakehira2025}, together with the degree argument in
\cite{WeiPeriod4}, show that the map $\beta(\alpha)=h(\alpha)$ supplies,
for $k\ge2$, a Galois-equivariant bijection between the root sets of
the two separable polynomials,
\[
 \beta:Z(\Gamma_k)\longrightarrow
 Z(\widetilde\Delta_{4k,4})
\]
for $\operatorname{Gal}(\overline{\mathbf Q}/\mathbf Q)$.  Restricting the
action to $\operatorname{Gal}(\overline{\mathbf Q}/B)$ preserves
equivariance.  The map $\beta$ therefore identifies the Galois orbits
over $B$, including their cardinalities.  Since the irreducible factors
of a separable polynomial correspond exactly to these orbits,
Theorem~\ref{thm:abelian-factorization} transfers without requiring an
explicit polynomial substitution:
\begin{theorem}\label{thm:delta-abelian-factorization}
For $k\ge2$ and every finite abelian extension $B/\mathbf Q$, the
period-$4$ delta factor $\widetilde\Delta_{4k,4}$ has exactly
\[
 [B\cap K_k:\mathbf Q]
\]
irreducible factors over $B$, each of degree
\[
 6[K_k:B\cap K_k].
\]
Over $\mathbf Q^{\mathrm{ab}}$ its complete factorization consists of
$\varphi(k)$ irreducible sextic factors.  Indeed, restricting $\beta$ to
$\operatorname{Gal}(\overline{\mathbf Q}/\mathbf Q^{\mathrm{ab}})$ still
identifies the relevant orbits.  More precisely, for every
primitive $k$-th root $\zeta$, the image under $\beta$ of
$Z(F_\zeta)$ is one orbit of cardinality six for
$\operatorname{Gal}(\overline{\mathbf Q}/\mathbf Q^{\mathrm{ab}})$.
The corresponding irreducible sextic delta factor is
\[
 H_\zeta(C)=\prod_{\alpha\in Z(F_\zeta)}
 \bigl(C-\beta(\alpha)\bigr)
 \in\mathbf Q^{\mathrm{ab}}[C].
\]
\end{theorem}

\begin{remark}
This is an orbit description; it does not identify the polynomial
$H_\zeta(C)$ in the parameter variable with $F_\zeta(T)$.
\end{remark}

For example, if $\gcd(k,\ell)=1$, then
$K_k\cap\mathbf Q(\zeta_\ell)=\mathbf Q$, and the period-$4$ delta
factor remains irreducible over $\mathbf Q(\zeta_\ell)$.

\section{The maximal abelian subfield of a parameter field}

\begin{corollary}[Maximal abelian subfield of a parameter field]
\label{cor:maximal-abelian-parameter-field}
Let $C$ be an exact-period-$4$ parabolic parameter with primitive
$k$-th root multiplier $\zeta$, where $k\ge2$.  Then
\[
 \mathbf Q(C)\cap\mathbf Q^{\mathrm{ab}}=\mathbf Q(\zeta).
\]
Thus $\mathbf Q(\zeta)$ is the largest subfield of $\mathbf Q(C)$ that
is Galois abelian over $\mathbf Q$, and $C\notin\mathbf Q^{\mathrm{ab}}$.
\end{corollary}

\begin{proof}
Let $\alpha$ be the corresponding root of $F_\zeta$, so
$C=\beta(\alpha)=h(\alpha)$.  The root correspondence is injective and
$\operatorname{Gal}(\overline{\mathbf Q}/\mathbf Q)$-equivariant.
Consequently, for every
$\sigma\in\operatorname{Gal}(\overline{\mathbf Q}/\mathbf Q)$,
\[
 \sigma(\alpha)=\alpha
 \quad\Longleftrightarrow\quad
 \sigma(C)=C.
\]
The two elements have the same stabilizer, and therefore
\[
 \mathbf Q(C)=\mathbf Q(\alpha).
\]
Since $\zeta=g(\alpha)$, the right-hand field already contains
$K=\mathbf Q(\zeta)$ and equals $K(\alpha)$.  The abelian-stability
theorem gives
\[
 K(\alpha)\cap K^{\mathrm{ab}}=K.
\]
Because $\mathbf Q^{\mathrm{ab}}/K$ is abelian,
$\mathbf Q^{\mathrm{ab}}\subset K^{\mathrm{ab}}$.  Moreover,
$K\subset\mathbf Q(C)$ and $K\subset\mathbf Q^{\mathrm{ab}}$.  The
upper inclusion from the abelian-stability theorem and this lower
inclusion give
\begin{equation}\label{eq:parameter-abelian-intersection}
 \mathbf Q(C)\cap\mathbf Q^{\mathrm{ab}}=\mathbf Q(\zeta).
\end{equation}
Thus the multiplier field is precisely the largest subfield
$L\subset\mathbf Q(C)$ for which $L/\mathbf Q$ is Galois abelian.  In
particular,
$C\notin\mathbf Q^{\mathrm{ab}}$: otherwise
$\mathbf Q(C)\subset\mathbf Q^{\mathrm{ab}}$ would force
$\mathbf Q(C)=K$, contradicting $[\mathbf Q(C):K]=6$.
\end{proof}

When the multiplier is $1$, the exact-period-$4$ parameters instead
satisfy
\[
 p(C)=C^3+9C^2+27C+135=0.
\]
After $U=C+3$, this becomes $U^3+108$.  A rational root of this monic
cubic would be an integer $r$ satisfying $r^3=-108$, which is
impossible because $-108$ is not an integer cube.  The cubic is
therefore irreducible.  Its discriminant is
\[
 -27\cdot108^2=-2^4 3^9,
\]
which is not a square.  Its Galois closure consequently has group
$S_3$, and the cubic field $\mathbf Q(C)$ is not Galois over $\mathbf Q$.  Every
finite subfield of $\mathbf Q^{\mathrm{ab}}$ is Galois over
$\mathbf Q$, so these parameters also lie outside
$\mathbf Q^{\mathrm{ab}}$.  Since a cubic field has no nontrivial proper
subfield, its largest subfield Galois abelian over $\mathbf Q$ is
$\mathbf Q$.  We conclude that no exact-period-$4$
parabolic parameter belongs to $\mathbf Q^{\mathrm{ab}}$.

\section{Torsion Hilbert irreducibility context}

For this interpretation one must use the finite cover of degree six
\[
 g:U\longrightarrow\mathbf G_m,
 \qquad
 U=\mathbf G_m\setminus g^{-1}(0).
\]
Its degree is the degree of the rational function $g$, namely six.
After writing the target coordinate as $X=g(T)$, the element $T$
satisfies the monic equation $P_X(T)=0$ with nonzero constant term;
both $T$ and $T^{-1}$ are integral over
$\overline{\mathbf Q}[X,X^{-1}]$, which also proves finiteness.  The
pullback condition means that
the fiber product with $x\mapsto x^n$ is geometrically irreducible for
every $n\ge1$.  Its function-field equation is
\[
 Y^n=g(T).
\]
Over $\overline{\mathbf Q}(T)$, the function $g$ is not a $p$-th power
for any prime $p$.  For $p\ge3$, this follows from its pole of order two
at $T=0$.  For $p=2$, suppose $g=w^2$.  Its only poles force
\[
 w=aT^2+bT+c+dT^{-1}.
\]
Comparison of the coefficients of $T^4,T^3,T^{-1},T^{-2}$ gives
\[
 a^2=-1,\qquad b=a,\qquad d^2=16,\qquad cd=4.
\]
Thus $c^2=1$, whereas comparison at $T^2$ gives
\[
 -1+2ac=-4,
 \qquad ac=-\frac32,
\]
so $(ac)^2=9/4$, contradicting $a^2c^2=-1$.  Hence $g$ is not a
square.

The field $\overline{\mathbf Q}(T)$ contains all roots of unity.  The
Kummer binomial criterion \cite{Lang2002} in this setting states that
$Y^n-q$ is irreducible if and only if $q$ is not a $p$-th power for
every prime $p\mid n$.  In the general binomial criterion one must also
exclude $q\in-4\overline{\mathbf Q}(T)^4$ when $4\mid n$; here every
element $-4b^4=(2ib^2)^2$ is already a square, so the preceding square
argument excludes this case as well.  The preceding valuation
and coefficient arguments therefore prove that $Y^n-g(T)$ is
irreducible for every $n\ge1$, which is precisely the pullback
condition.

Zannier's cyclotomic Hilbert irreducibility theorem
\cite[Theorem 2.1]{Zannier2010} says that a cover defined over the cyclotomic
closure and satisfying this condition has irreducible fibers over that
field outside a finite union of proper torsion cosets.  For
$\mathbf G_m$, such a union is a finite set of torsion points.  Thus the
theorem gives finitely many exceptional torsion fibers.  By the
Kronecker--Weber theorem \cite{Washington1997}, the cyclotomic closure of
$\mathbf Q$ is $\mathbf Q^{\mathrm{ab}}$.  The abelian-stability
theorem determines them exactly: for $\zeta\ne1$, the inclusion
$\mathbf Q^{\mathrm{ab}}\subset\mathbf Q(\zeta)^{\mathrm{ab}}$ shows
that $F_\zeta$ is irreducible over $\mathbf Q^{\mathrm{ab}}$, whereas
\[
 F_1(T)=(T+2)(T^2+4)(T^3-2).
\]
Thus $\zeta=1$ is the unique exceptional torsion fiber for this cover.
